\section{The Data Kernel Perspective Space}

Following \cite{aranyak}, we consider a generative model $ f $ with weights $ W $ and decoding function (temperature, maximum response length, etc.) $ \delta $ to be a random map from a query space $ \mathscr{Q} $ to a response space $ \mathcal{X} $.
We let $ f_{1}, \hdots, f_{n} $ be $ n $ such random maps for a shared $ \mathscr{Q} $ and $\mathcal{X} $. 
For example, $ f_{i} $ may be the neural model parameterized by the fixed context augmentation $ a_{i} $, the retrieval vector database $ V_{i} $, the fine-tuned model with weights $ W + \Delta W_{i} $ or the model with decoding function $ \delta_{i} $. 
Given a collection of queries $ Q = \{q_{1}, \hdots, q_{m}\} $ 
% assumed to be $ i.i.d. $ realizations from a query distribution $ G $ with support $ \mathscr{Q} $, 
we let $ F_{ij} $ denote the distribution on the response space for fixed $ f_{i} $ and $ q_{j} $ and assume model responses $ f_{i}(q_{j})_{1}, \hdots, f_{i}(q_{j})_{r} = x_{ij1}, \hdots, x_{ijr} $ are $ i.i.d. $ realizations from $ F_{ij} $. 
% We assume the query space $Q$ and the embedded response space $g(\mathcal{X})$ to be bounded. 

Our goal is to induce a representation of the neural models that captures information relevant to model-level statistical inference.
We abuse notation by letting $ x $ represent a realization of a random variable and the random variable itself and include necessary context when there is potential for confusion. 
% -- such as predicting if a model has had access to sensitive information or predicting a model's performance on a benchmark. 

% \subsection*{Methodology}
We let $ g: \mathcal{X} \to \mathbb{R}^{p} $ be an embedding function that maps a response to a $p$-dimensional real-valued vector and, for notational convenience, let $ \bar{x}_{ij} = \frac{1}{r} \sum_{k=1}^{r} g(x_{ijk}) $ be the average of the embedded responses. 
We use $ \bar{X}_{i} \in \mathbb{R}^{m \times p}$ to denote the matrix where the $ j $th row is $ \bar{x}_{ij} $ and view this matrix as a representation of model $ f_{i} $ with respect to $ Q $ and $ g $. 

We define $ D $ as the $ n \times n $ pairwise distance matrix with entries 
\begin{align}
\label{euc-dist-matrix}
    D_{ii'} = \frac{1}{m} \big|\big| \bar{X}_{i} - \bar{X}_{i'} \big|\big|_{F}.
\end{align} Each entry $ D_{ii'} $ is the average difference between the average embedded response between model $ f_{i} $ and model $ f_{i'} $ and captures the difference in average model behavior with respect to $ Q $ and $ g $. 
With the form described by Eq. \eqref{euc-dist-matrix}, $ D $ is a Euclidean distance matrix and the multidimensional scaling (MDS) of $ D $ yields $ d$-dimensional Euclidean representations of the matrices $ \bar{X}_{i} $ \citep{torgerson1952multidimensional} and, thus, Euclidean representations of the models $ f_{1}, \hdots, f_{n} $ with respect to $ Q $ and $ g $.
Letting $ \widehat{\psi} := \text{MDS}(D) $, we refer to the $ i $th row of $ \widehat{\psi}_{i} \in \mathbb{R}^{d} $ as the \textit{perspective} of model $ f_{i} $ and the entire configuration $ \widehat{\psi} $ as the \textit{data kernel perspective space} (DKPS).



% In particular, let $ D^{2} $ be the $ n \times n $ matrix with entries $ D_{ii'}^{2} $, $ e_{n} \in \mathbb{R}^{n}$ be the vector of all ones, $ I_{n} $ be the $ n \times n $ identity matrix, and $ P = I_{n} - e_{n}e_{n}^{T} / n $. Then, letting $ B = -\frac{1}{2} P D^{2} P^{T} $ and $ v_{\ell} $ be the eigenvector associated with the $ \ell $th largest eigenvalue $ \lambda_{\ell} $ of $ B $, the rows of the matrix $ \widehat{\psi} \in \mathbb{R}^{n \times d} $ with columns $ \lambda_{\ell}^{2} v_{\ell} $ are $ d $-dimensional Euclidean representations of the models with respect to $ Q $. 



% \subsubsection*{Properties of the Euclidean representations}

% \clearpage
\subsection{Analytical properties of the data kernel perspective space}

While the perspectives are Euclidean objects, it is not possible to comment on their properties without imposing constraints on the $ F_{ij} $.
To facilitate analysis we assume each query is from the query distribution $ G $ with support $ \mathscr{Q} $  
% let the distribution of response distributions $ F_{ij} $ with $ q_{j} \sim G $ be $ \mathcal{F}_{i}(G) $ for fixed $ f_{i} $ (i.e., $ F_{ij} \sim \mathcal{F}_{i}(G) $) 
and let $ \mu_{i} \in \mathbb{R}^{m \times p} $ denote the matrix whose $ i^{th}$ row is $ \mathbb{E}_{x \sim F_{ij}}\left[g(x)\right]$. 
We let $ \Delta $ be such that $ \Delta_{ii'} = \frac{1}{m}||\mu_{i} - \mu_{i'}||_{F} $.
For settings where $ m $ grows and $ q_{j} \iid G $, we refer to $ \mu_{i} $ as $ \mu_{i}(G) $ and let $ \Delta^{*}(G) $ be such that $ \Delta^{*}_{i i'}(G) = \lim_{m \to \infty} \frac{1}{m} || \mu_{i} - \mu_{i'}||_{F} $.
% \begin{align*}
% \mu_{i}(G)~:=~\mathbb{E}_{F_{ij}\sim\mathcal{F}(G)}\left[\mathbb{E}_{x\sim F_{ij}}\left[g(x))\right]\right].
% \end{align*}
Finally, we assume $ g $ is bounded.\footnote{Boundedness is typically satisfied in practice for well-defined $ g $: for language models, an element of $ \mathcal{X} $ is a finite sequence from a finite vocabulary; for text-to-image models $ \mathcal{X} = \{0, \hdots, 255\}^{3}$ is finite.

% . For example, for language models elements of both $ \mathscr{Q} $ and $ \mathcal{X} $ are sequences of tokens from a finite vocabulary where the length of an element in $ \mathscr{Q} $ is bounded above by the size of the context window and the length of an element of $ \mc{X} $ is bounded above by a user-defined maximum sequence length.
} 
%( i.e., $ |\mathscr{Q}| < \infty $ and $ |\mathcal{X}| < \infty $).  
% Together, these imply that the image of $ \mathcal{X} $ under $ g $ is bounded and, hence, that $ g(x) $ with $ x \sim F_{ij} $ has finite moments and $ \mathbb{E}_{F_{ij} \sim \mathcal{F}(G)}\left[\mathbb{E}_{x \sim F_{ij}}\left[g(x)\right] \right] $ is finite. 

With $ q_{1}, \hdots, q_{m} \iid G $, and $ x_{ij1}, \hdots, x_{ijr} \iid F_{ij} $, 
% and letting 
% \begin{align*}
% \mu_{i}(G)~:=~\mathbb{E}_{F_{ij}\sim\mathcal{F}(G)}\left[\mathbb{E}_{x\sim F_{ij}}\left[g(x))\right]\right],
% \end{align*}
we have
% $|D_{i i'}-\Delta_{i i'}(G)| \to 0$  as $m,r \to \infty$ 
% with high probability for all $ i, i'$.
% where $\Delta_{i i'}(G)=\frac{1}{m}
% \left\lVert 
% \mu_i(G)-\mu_{i'}(G)
% \right\rVert
% $. 
% Additionally, we assume 
% $\lim_{m,r \to \infty}
% \Delta_{i i'}(G)
% \to \Delta^{(\infty)}(G)
% $  for all $n$, where $\Delta^{(\infty)}
% \in \mathbb{R}^{n \times n}
% $ is a dissimilarity matrix, so we have
% for all $n$,
\begin{equation}
    D_{i i'}=
    \frac{1}{m}
    \left\lVert 
\bar{X}_i-\bar{X}_{i'}
    \right\rVert
    \to 
    \Delta^{*}_{i i'}(G)
\end{equation}
with high probability
as $m,r \to \infty$ for all $ (i, i') \in \{1, \hdots, n\} \times \{1, \hdots, n\} $.
%\begin{align}
%\label{eq:entry-of-D2}
%    D_{ii'} = \frac{1}{m} \big\| \big(\bar{X}_{i} - \bar{X}_{i'}\big) \big\|_{F} 
%    \xrightarrow[]{}
%    \big\|  \mu_{i}(G) - \mu_{i'}(G) \big\|_{F} =: \Delta_{ii'}(G)
%\end{align} 
%as $ R, m \to \infty $ for all $ (i, i') \in \{1, \hdots, n\} \times \{1, \hdots, n\} $ \citep{aranyak}.  
Under technical assumptions described in \cite{aranyak}, there exists $ \psi(G) := \text{MDS}(\Delta^{*}(G)) $ such that $ \widehat{\psi} \to \psi(G) $. The configuration $ \psi(G) $ captures the true geometry of the mean descrepancies of the model responses with respect to queries from $ G $.
We drop the argument of $ \psi $ when the distribution of the queries is clear.
% and hinges on the boundedness of $ g(x) $. 
In a more general setting where the models are assumed to be $ i.i.d. $ realizations from a model distribution $ F_{model}$, the distance matrix $ D $ converges to $ \Delta^{*} $ and, under technical assumptions, there exists $ \psi $ such that $ \widehat{\psi} \to \psi $ as $ m, r \to \infty $, for all $n$ \citep{aranyak}.


In settings where the query space contains only a single element or where each $ F_{ij} $ is a point mass on $ x_{ij} $, $ \Delta_{ii'} $ is exactly the maximum mean discrepancy \citep{gretton2012kernel} of the distributions $ F_{ij} $ and $ F_{i'j} $ under a linear map applied to $ g(x) $. 
For settings where $ | \mathscr{Q} | > 1 $ and the $ F_{ij} $ are not all point masses, the analysis is more complicated since both $ F_{ij} $ and $ \bar{x}_{ij}$ 
are random variables. 
% The entry-wise convergence of $ D $ to $ \Delta(G) $ implies the convergence of $ Z_{i} $ to an analogously defined $ Z^{*}_{i}(G) \in \mathbb{R}^{p} $ up to an orthogonal transformation. The proof of this is provided in Appendix [ref].


% While we do not explore alternative dissimilarities on the $ m \times p \times R $ matrices obtained after embedding the $ R $ samples from each $ F_{ij} $ in this paper, we note that it is possible (and perhaps worthwhile) to consider more expressive notion of dissimilarity on the models. 
% For example, we could replace the row-wise difference with the average of a row-wise MMD with a universal kernel to ensure that we capture all relevant differences between the response distributions induced by different $ f_{i} $.
% The resulting $ n \times n $ pairwise distance matrix will likely not be a Euclidean distance matrix, though transforming it into useful and analyzable representations of the models is sometimes possible [cite, Trosset?].

% \clearpage

\subsection{Statistical inference in the data kernel perspective space}
\label{subsec:inference}

Consider the classical statistical learning problem \citep[Chapter~2]{hastie2009elements}:
Given training data $ \mathcal{T}_{n} = \{ (X_{1}, y_{1}), \hdots, (X_{n}, y_{n}) \} $ assumed to be  $ i.i.d. $ realizations from the joint distribution $ P_{XY} $, select the decision function $ h: \mathcal{X} \to \mathcal{Y} $ that minimizes the expected value of a loss function $ \ell $ with respect to $ P_{XY} $ within a class of decision functions $ \mathcal{H} $ for a test observation $ X $ assumed to be an independent realization from the marginal distribution $ P_{X}$. Or, with $ \mathcal{R}_{\ell}(P_{XY}, h) := \mathbb{E}_{P_{XY}}\left[\ell(h(X), y)\right] $, select $ h $ such that \begin{align*}
\label{eq:risk}
    h \in \argmin_{h \in \mathcal{H}} \mathcal{R}_{\ell}(P_{XY}, h). 
    % \min_{h \in \mathcal{H}} \mathbb{E}_{\mathcal{P}_{XY}}\left[\ell(h(X), y)\right]
    % = \min_{h \in \mathcal{H}} \mathcal{R}(\mathcal{P}_{XY}, \mathcal{H}).
\end{align*}
We let $ \mathcal{R}_{\ell}^{*}(P_{XY}, \mathcal{H}) $ denote the expected loss of $ h $ in the argmin and let $ \mathcal{R}^{*}_{\ell}(P_{XY}) $ denote the minimum $ \mathcal{R}_{\ell}^{*}(P_{XY}, \mathcal{H}) $ over all $ \mathcal{H} $.

The joint distribution $ P_{XY} $ is often unavailable and the decision function is selected based on $ \mathcal{T}_{n} $. 
We let $ h(\; \cdot \;; \mathcal{T}_{n}) $ denote such a decision function and say \textit{the sequence of decision functions $ \left(h(\; \cdot \;; \mathcal{T}_{1}), \hdots, h(\;\cdot\;; \mathcal{T}_{n})\right) $ is consistent for $ P_{XY} $ with respect to $ \mathcal{H}$} if \begin{align*}
Pr\big(\;\big|  \mathcal{R}_{\ell}(P_{XY}, h(\;\cdot\;; \mathcal{T}_{n})) - \mathcal{R}^{*}_{\ell}(P_{XY}, \mathcal{H}) \big| > \epsilon\big) \to 0
\end{align*} as $ n \to \infty $.

In our setting we observe $ \widehat{\mathcal{T}}_{n} = \{(\widehat{\psi}_{1}, y_{1}), \hdots, (\widehat{\psi}_{n}, y_{n})\} $ and the test observation $ \widehat{\psi}_{n+1} $ where $ \widehat{\psi}_{i} $ is an errorful observation of the true-but-unknown $ \psi_{i} $ and $ y_{i} $ is the model-level covariate corresponding to $ f_{i} $. 
% Letting $ \{(f_{1}, y_{1}), \hdots, (f_{n}, y_{n})\} $ be $ i.i.d. $ realizations from a distribution on generative models and $ P_{\psi Y} $ be the corresponding distribution induced by MDS($\Delta$),
Two important extensions of the convergence results described in \citep{aranyak} are the following theorems:
\begin{theorem}
\label{thm:rule-convergence}
    Under technical assumptions described in Appendix \ref{app:proofs}, \begin{align*}
    \mathcal{R}_{\ell}(P_{\psi Y}, h(\; \cdot \;; \widehat{\mathcal{T}}_{n})) \to \mathcal{R}_{\ell}(P_{\psi Y}, h(\; \cdot \;; \mathcal{T}_{n})) 
    \end{align*} as $ m, r \to \infty $;
\end{theorem}
% \begin{align}
% \label{eq:rule-convergence}
%     \mathcal{R}_{\ell}(P_{\psi Y}, h(\; \cdot \;; \widehat{\mathcal{T}}_{n})) \to \mathcal{R}_{\ell}(P_{\psi Y}, h(\; \cdot \;; \mathcal{T}_{n}))  
% \end{align}
and
\begin{theorem}
\label{thm:consistency}
    Under technical assumptions described in Appendix \ref{app:proofs}, if $ (h(\; \cdot \;; \mathcal{T}_{1}), \hdots, h(\; \cdot \;; \mathcal{T}_{n})) $ is consistent for $ P_{\psi Y} $ with respect to
    $ \mathcal{H} $ then $ (h(\; \cdot \;; \widehat{\mathcal{T}}_{1}), \hdots, h(\; \cdot \;; \widehat{\mathcal{T}}_{n})) $ is consistent for 
    $ P_{\psi Y} $ 
    with respect to
    $ \mathcal{H} $ as $ n, m, r \to \infty $.
\end{theorem}
% \begin{align}
% \label{eq:consistency}
%     &\text{if } (h(\; \cdot \;; \mathcal{T}_{1}), \hdots, h(\; \cdot \;; \mathcal{T}_{n})) \text{ is consistent for } P_{\psi Y} 
%     \text{ with respect to }
%     \mathcal{H}  \\ 
%     \notag &\text{then } 
%     (h(\; \cdot \;; \widehat{\mathcal{T}}_{1}), \hdots, h(\; \cdot \;; \widehat{\mathcal{T}}_{n})) 
%     \text{ is consistent for }
%     P_{\psi Y} 
%     \text{ with respect to }
%     \mathcal{H}
% \end{align}
% as $ n, m $ and $ R \to \infty $.
% \begin{enumerate}[label=\roman*.]
%     \item  $ \mathcal{R}_{\ell}(P_{\psi Y}, h(\; \cdot \;; \widehat{\mathcal{T}}_{n})) \to \mathcal{R}_{\ell}(P_{\psi Y}, h(\; \cdot\;, \;; \mathcal{T}_{n})) $ as $ m, R \to \infty $; and
%     \item if $ (h(\;\cdot
%     \;; \mathcal{T}_{1}), \hdots, h(\;\cdot\;; \mathcal{T}_{n})) $ is consistent for $ P_{
%     \psi Y} $ with respect to $ \mathcal{H} $ \\ then $ (h(\; \cdot \;; \widehat{\mathcal{T}}_{1}), \hdots, h(\; \cdot\;; \widehat{\mathcal{T}}_{n})) $ is consistent for $ P_{\psi Y} $ with respect to $ \mathcal{H} $ \\ as $ n, m $ and $ R \to \infty $.
% \end{enumerate}
The proofs of Theorems \ref{thm:rule-convergence} and \ref{thm:consistency} are provided in Appendix \ref{app:proofs}. 

These results do not say that with enough models, queries, and responses that inference using the $ \widehat{\psi} $ is ``optimal". 
Instead, Theorem \ref{thm:rule-convergence} says that with enough queries and responses for each query that the performance of a decision function that uses $ \widehat{\mathcal{T}}_{n} $ will be ``close" to the performance of the decision function that uses $ \mathcal{T}_{n} $ -- regardless of how far away $ \mathcal{R}_{\ell}(P_{\psi Y}, h(\; \cdot \;; \mathcal{T}_{n})) $ is from $ \mathcal{R}_{\ell}^{*}(P_{\psi Y}, \mathcal{H}) $.

Similarly, Theorem \ref{thm:consistency} says that if there is a sequence of decision functions known to be consistent with respect to $ \mathcal{H} $ when using $ \mathcal{T}_{n} $, then the same sequence of decision functions is consistent with respect to $ \mathcal{H} $ when using $ \widehat{\mathcal{T}}_{n} $ where $ m $ and $ r $ are allowed to grow. 
The result does not provide instructions on how to construct such a sequence. Nor does it provide insight as to how close $ \mathcal{R}_{\ell}^{*}(P_{\psi Y}, \mathcal{H}) $ is to $ \mathcal{R}_{\ell}^{*}(P_{\psi Y}) $. 
Most importantly -- and most unique to our setting -- Theorem \ref{thm:consistency} does not provide guidance for how to select a query distribution $ G $ such that $ | \mathcal{R}_{\ell}^{*}(P_{f Y}) - \mathcal{R}_{\ell}^{*}(P_{\psi Y}) | $ is minimal. 

While we do not provide theoretical details as to which situations optimal model performance is possible, if $ y_{i} $ can be described as a function of the average of model responses for a set of queries from an optimal $ G^{*} $ then $ \psi(G^{*}) $ will be such that $ \mathcal{R}_{\ell}^{*}(P_{f Y}) = \mathcal{R}_{\ell}^{*}(P_{\psi Y}) $ \citep[Chapter~32]{devroye2013probabilistic}.
% Common model-level covariates such as the performance on a benchmark, a model's toxicity and bias, or whether a model is likely to produce sensitive information can be described this way.
Conversely, it is unclear how effective decision functions constructed with  $ \psi(G^{*}) $ for model-level covariates that cannot be described as a function of the average response will behave. 
Further, it is unclear how ``close" a query distribution $ G $ needs to be to $ G^{*} $ to induce an informative representation of the models with respect to $ y_{i} $. 
We demonstrate a potential trade-off between the proximity of $ G $ to $ G^{*} $ and the number of queries necessary to tease out model differences in one of the experimental settings below. 

Lastly, we note that our setting, methodology, and theorems are general to classes of generative models in which there exists a well-defined embedding function from the output space to $ \mathbb{R}^{p} $, though our experimental settings focus entirely on collections of language models.

\subsection{An illustrative example -- ``Was RA Fisher great?"}
\label{subsec:fisher}
\begin{figure*}[t!]
    \centering
    \begin{subfigure}[t]{0.49\textwidth}
        \centering
        \includegraphics[width=\textwidth]{figures/was-ra-fisher-great/3d-surface-plot-091524.png}
    \end{subfigure}
    ~ 
    \begin{subfigure}[t]{0.49\textwidth}
        \centering
        \includegraphics[width=\textwidth]{figures/was-ra-fisher-great/ideal_fig1_nmr.png}
    \end{subfigure}
    \caption{
    \textbf{Left.} The 2-d Data Kernel Perspective Space (DKPS) and covariate surface for a collection of 550 models parameterized by fixed augmentations.
    \textbf{Right.} The performance of the 1-nearest neighbor regressor in DKPS for predicting the probability that an unlabeled model responds ``yes" to ``Was RA Fisher great?". 
    % The first dimension of DKPS correlates strongly with the covariate. Performance of the 1-NN regressor depends on both the number of models with known covariates and the number of queries used to induce the DKPS.
    }
    \label{fig:ra-fisher}
\end{figure*}

The first model-level inference task we study is a toy example where the task is to predict the probability that a model will respond ``yes" to the question ``Was RA Fisher great?". 
The question is chosen due to its subjectiveness -- there is no correct answer to the question of a person's greatness -- and its duality -- Ronald A. Fisher is considered one of the most influential statisticians in history and is considered an advocate of the 20th century Eugenics movement \citep{bodmer2021outstanding}. 

We use a 4-bit version of Meta's \texttt{LLaMA-2-7B-Chat} \citep{touvron2023llama} as a base model and consider a collection of models parameterized by fixed context augmentations. 
Each augmentation $ a_{i} $ contains information related to RA Fisher's statistical achievements (i.e., $ a_{i} $ = ``RA Fisher pioneered the principles of the design of experiments") or to his involvement in the eugenics movement (i.e., $ a_{i} $ = ``RA Fisher's view on eugenics were primarily based on anecdotes and prejudice.") and is prepended to every query.
The covariate corresponding to a given model is calculated by prompting the base model with the appropriately formatted prompt ``Give a precise answer to the question based on the context. Don't be verbose. The answer should be either a yes or a no. $ a_{i} $. Was RA Fisher great?" until there are 100 valid responses. We let $ y_{i} $ be the average number of ``yes"es.

To induce a DKPS for this task, we consider queries sampled from OpenAI's ChatGPT with the prompt ``Provide 100 questions related to RA Fisher.".  
For a given query $ q_{j} $ we prompt the base model with the appropriately formatted prompt ``$a_{i} \; q_{j}$" and fix $ r = 1 $.
The left figure of Figure \ref{fig:ra-fisher} is a 3-d figure where the first two dimensions are the DKPS of the $ n = 550 $ ($ 275 $ statistics augmentations, $ 275 $ eugenics augmentations) models induced with $ m = 100 $ queries. Model responses are embedded by averaging the per-token last layer activation of the base model.  
The third dimension is an interpolated $ y_{i} $ surface with a linear kernel \citep{du2008radial}. 
The first dimension of the DKPS is clearly capable of distinguishing between models adorned with ``statistics" augmentations and models adorned with ``eugenics" augmentations. 
The shape of the interpolated covariate surface is highly correlated with this feature.
A description of the augmentations that parameterize the models and the queries used to induce the DKPS is provided in Appendix \ref{app:fisher}.

The right figure of Figure \ref{fig:ra-fisher} shows the performance of a 1-nearest neighbor (1-NN) regressor in DKPS for a varying number of labeled models and a varying number of queries. 
The DKPS is induced with $ n $ models and the regressor is trained with $ n - 1 $ of the model-level covariates. 
The mean squared error reported is the error of the 1-NN regressor for predicting the ``left out" model's covariate ($\pm$ 1 S.E.).
As Theorems \ref{thm:rule-convergence} and \ref{thm:consistency} suggest, the performance of the regressor is dependent on both the number of models and the number of queries: the more models and the more queries the better.
The scale of the impact of more models and more queries, however, depends on its counterpart. 
The amount of models does not have a large impact on predictive performance if the number of queries is small. 
The amount of queries has a large impact on performance regardless of the number of models.